\documentclass[ignorenonframetext, 11pt, table]{beamer}
\gdef\pathtomaster{../_LaTeX_master}
%-- Nummerierung bis Subsubsection für Report
\setcounter{secnumdepth}{3}	% zählt auch subsubsection
\setcounter{tocdepth}{3}	% Inhaltsverzeichnis bis in subsubsection

% LuaLaTeX Unterscheidung
\usepackage{ifluatex}

% Input inkl. Umlaute, Silbentrennung
\ifluatex
	\usepackage{fontspec}
\else
	\usepackage[utf8]{inputenc}
	\usepackage[T1]{fontenc}
\fi
\usepackage[english,ngerman]{babel}
\usepackage{csquotes}	% Anführungszeichen
\usepackage{marvosym}
\usepackage{eurosym}

% Style-Aufhübschung
\usepackage{soul, color}	% Kapitälchen, Unterstrichen, Durchgestrichen usw. im Text
%\usepackage{titleref}
\usepackage[breakwords, fit]{truncate}	% Abschneiden von Sätzen
\renewcommand{\TruncateMarker}{\,…}

% Mathe usw.
\usepackage{amssymb}
\usepackage{amsthm}
\usepackage[fleqn,intlimits]{amsmath}	% fleqn: align-Umgebung rechtsbündig; intlimits: Integralgrenzen immer ober-/unterhalb
%\usepackage{mathtools} % u.a. schönere underbraces
\usepackage{xcolor}
\usepackage{esint}	% Schönere Integrale, \oiint vorhanden
\everymath=\expandafter{\the\everymath\displaystyle}	% Mathe Inhalte werden weniger verkleinert
\usepackage{wasysym}	% mehr Symbole, bspw \lightning
%\renewcommand{\int}{\int\limits}
%\usepackage{xfrac}	% mehr fracs: sfrac{}{}
\let\oldemptyset\emptyset	% schöneres emptyset
\let\emptyset\varnothing
%\RequirePackage{mathabx}	% mehr Symbole
\mathchardef\mhyphen="2D	% Hyphen in Math
\usepackage{mathtools}		% zum Beispiel {pmatrix}[r] für Formatierung der Matrix

% tikz usw.
\usepackage{tikz}
\usepackage{pgfplots}
\pgfplotsset{compat=1.11}	% Umgeht Fehlermeldung
\usetikzlibrary{graphs}
%\usetikzlibrary{through}	% ???
\usetikzlibrary{arrows}
\usetikzlibrary{arrows.meta}	% Pfeile verändern / vergrößern: \draw[-{>[scale=1.5]}] (-3,5) -> (-3,3);
\usetikzlibrary{automata,positioning} % Zeilenumbruch im Node node[align=center] {Text\\nächste Zeile} automata für Graphen
\usetikzlibrary{matrix}
\usetikzlibrary{patterns}	% Schraffierte Füllung
\usetikzlibrary{shapes.geometric}	% Polygon usw.
\tikzstyle{reverseclip}=[insert path={	% Inverser Clip \clip
	(current page.north east) --
	(current page.south east) --
	(current page.south west) --
	(current page.north west) --
	(current page.north east)}
% Nutzen: 
%\begin{tikzpicture}[remember picture]
%\begin{scope}
%\begin{pgfinterruptboundingbox}
%\draw [clip] DIE FLÄCHE, IN DER OBJEKT NICHT ERSCHEINEN SOLL [reverseclip];
%\end{pgfinterruptboundingbox}
%\draw DAS OBJEKT;
%\end{scope}
%\end{tikzpicture}
]	% Achtung: dafür muss doppelt kompliert werden!
\usepackage{graphpap}	% Grid für Graphen
\tikzset{every state/.style={inner sep=2pt, minimum size=2em}}
\usetikzlibrary{mindmap, backgrounds}
%\usepackage{tikz-uml}	% braucht Dateien: http://perso.ensta-paristech.fr/~kielbasi/tikzuml/

% Tabular
\usepackage{longtable}	% Große Tabellen über mehrere Seiten
\usepackage{multirow}	% Multirow/-column: \multirow{2[Anzahl der Zeilen]}{*[Format]}{Test[Inhalt]} oder \multicolumn{7[Anzahl der Reihen]}{|c|[Format]}{Test2[Inhalt]}
\renewcommand{\arraystretch}{1.3} % Tabellenlinien nicht zu dicht
\usepackage{colortbl}
\arrayrulecolor{gray}	% heller Tabellenlinien
\usepackage{array}	% für folgende 3 Zeilen (für Spalten fester breite mit entsprechender Ausrichtung):
\newcolumntype{L}[1]{>{\raggedright\let\newline\\\arraybackslash\hspace{0pt}}m{\dimexpr#1\columnwidth-2\tabcolsep-1.5\arrayrulewidth}}
\newcolumntype{C}[1]{>{\centering\let\newline\\\arraybackslash\hspace{0pt}}m{\dimexpr#1\columnwidth-2\tabcolsep-1.5\arrayrulewidth}}
\newcolumntype{R}[1]{>{\raggedleft\let\newline\\\arraybackslash\hspace{0pt}}m{\dimexpr#1\columnwidth-2\tabcolsep-1.5\arrayrulewidth}}
\usepackage{booktabs}	% Kompatibilität für array und \hrule
\usepackage{caption}	% Um auch unbeschriftete Captions mit \caption* zu machen

% Nützliches
\usepackage{verbatim}	% u.a. zum auskommentieren via \begin{comment} \end{comment}
\usepackage{tabto}	% Tabs: /tab zum nächsten Tab oder /tabto{.5 \CurrentLineWidth} zur Stelle in der Linie
\NumTabs{6}	% Anzahl von Tabs pro Zeile zum springen
\usepackage{listings} % Source-Code mit Tabs
\usepackage{lstautogobble} 
\usepackage{enumitem}	% Anpassung der enumerates
%\setlist[enumerate,1]{label=(\arabic*)}	% global andere Enum-Items
\renewcommand{\labelitemiii}{$\scriptscriptstyle ^\blacklozenge$} % global andere 3. Item-Aufzählungszeichen
\usepackage{letltxmacro} % neue Definiton von Grundbefehlen
% Nutzen:
%\LetLtxMacro{\oldemph}{\emph}
%\renewcommand{\emph}[1]{\oldemph{#1}}
\RequirePackage{xpatch}	% ua. Konkatenieren von Strings/Variablen (etoolbox)
\usepackage{xstring}	% String Operationen
\usepackage{minibox}	% Minibox anstatt \fbox{} für Boxen mit (manuellen) Zeilenumbruch
\usepackage{framed}		% \begin{framed} statt minibox/fbox mit automatischem Zeilenumbruch
\usepackage{microtype}	% Verfeinerte Hyphenation
\usepackage{import}		
\usepackage{standalone}	% Zusammen mit import: Einfügen von separaten Latex-Dateien (die schon header haben)

% Einrichtung von lst
%\renewcommand*{\ttdefault}{qcr}
\lstset{
basicstyle=\ttfamily, 
%basicstyle=\footnotesize\sffamily\color{black},
%mathescape=true, 
%escapeinside=^^, 
autogobble, 
tabsize=2,
frame=single,
rulecolor=\color{lightgray},
numbers=left,
numbersep=5pt,
numberstyle=\tiny\color{gray},
commentstyle=\color{gray},
keywordstyle=\color{htwblue},
stringstyle=\color{htworange},
morecomment=[l][\color{magenta}]{\#}
showspaces=false,
showstringspaces=false,
breaklines=true,
literate=%
    {Ö}{{\"O}}1
    {Ä}{{\"A}}1
    {Ü}{{\"U}}1
    {ß}{{\ss}}1
    {ü}{{\"u}}1
    {ä}{{\"a}}1
    {ö}{{\"o}}1
    {„}{{\glqq}}1
    {“}{{\grqq}}1
    {~}{{\textasciitilde}}1
}
\usepackage{scrhack} % Fehler umgehen
\def\ContinueLineNumber{\lstset{firstnumber=last}} % vor lstlisting. Zum wechsel zum nicht-kontinuierlichen muss wieder \StartLineAt1 eingegeben werden
\def\StartLineAt#1{\lstset{firstnumber=#1}} % vor lstlisting \StartLineAt30 eingeben, um bei Zeile 30 zu starten
\let\numberLineAt\StartLineAt

% BibTeX
\usepackage[
backend=bibtex,
bibencoding=ascii,
style=authoryear, citestyle=authoryear-ibid,
%doi=false,
%isbn=false,
%url=false
]{biblatex}	% BibTeX
\usepackage{makeidx}
%\makeglossary
%\makeindex
\usepackage{url}	% Darstellung von URLs

% Grafiken
\usepackage{graphicx}
\usepackage{epstopdf}	% eps-Vektorgrafiken einfügen
\usepackage{transparent}	% transparent nutzen: {\transparent{0.4} ...}
%\epstopdfsetup{outdir=\customDir}
% Prüft, ob Grafik existiert (mit \ifvalidimage{}{}) [Quelle: https://tex.stackexchange.com/a/99176]:
\makeatletter
\newif\ifgraphicexist
\catcode`\*=11
\newcommand\ifvalidimage[1]{%
    \begingroup
    \global\graphicexisttrue
    \let\input@path\Ginput@path
    \filename@parse{#1}%
    \ifx\filename@ext\relax
    \@for\Gin@temp:=\Gin@extensions\do{%
        \ifx\Gin@ext\relax
        \Gin@getbase\Gin@temp
        \fi}%
    \else
    \Gin@getbase{\Gin@sepdefault\filename@ext}%
    \ifx\Gin@ext\relax
    \global\graphicexistfalse
    \def\Gin@base{\filename@area\filename@base}%
    \edef\Gin@ext{\Gin@sepdefault\filename@ext}%
    \fi
    \fi
    \ifx\Gin@ext\relax
    \global\graphicexistfalse
    \else 
    \@ifundefined{Gin@rule@\Gin@ext}%
    {\global\graphicexistfalse}%
    {}%
    \fi  
    \ifx\Gin@ext\relax 
    \gdef\imageextension{unknown}%
    \else
    \xdef\imageextension{\Gin@ext}%
    \fi 
    \endgroup 
    \ifgraphicexist
    \expandafter \@firstoftwo
    \else
    \expandafter \@secondoftwo
    \fi 
} 
\catcode`\*=12
\makeatother
\LetLtxMacro{\forceincludegraphics}{\includegraphics}	% neuer Befehl für includegraphics
\renewcommand{\includegraphics}[2][]{	% altes includegraphics neu definieren, damit es auch nicht vorhandene einfügt
\ifvalidimage{#2}{
\forceincludegraphics[#1]{#2}
}{
\message{Achtung: Grafik wurde nicht gefunden: '#2'}
\minibox[frame]{
\textbf{\StrSubstitute{#2}{_}{\_}}%  \ifthenelse{\isempty{#1}}{}{\\\textit{#1}}
}}}

% pdf-Setup
\usepackage{pdfpages}
\usepackage[bookmarks,%
bookmarksopen=false,% Klappt die Bookmarks in Acrobat aus
colorlinks=true,%
linkcolor=black,%
citecolor=black,%
urlcolor=htworange,%
]{hyperref}

% Nach Dokumentenbeginn ausführen:
\makeatletter
\AtBeginDocument{	
	% pdf Informationen
	\hypersetup{
		pdfauthor={\@author},
		pdftitle={\@title},
		pdfproducer={LaTeX Layout erstellt von Falk-Jonatan Strube}
	}
	% signierte Fußzeile
	%\providecommand{\ofootname}{Mitschrift von}
	\providecommand{\pathtomaster}{../}
	\ofoot{%
	%\ifdefined\ofoottext
	%\else
	%\begin{minipage}{40mm}
	%\begin{flushright}
	%\textcolor{gray}{
	%\ofootname\\
	%Falk-Jonatan Strube\vspace*{-1mm}}
	%\end{flushright}
	%\end{minipage}
	%\fi%
	%\hspace*{1mm}%
	\begin{minipage}{10 mm}
	\includegraphics[height=9mm]{\pathtomaster/FJS-logo_gray}
	\end{minipage}%
	}
}
\makeatother



%%% ==================================================
%%% ===================== MACROS =====================
%%% ==================================================
% Horizontale Linie:
\newcommand{\HRule}[1][\medskipamount]{\par
  \vspace*{\dimexpr-\parskip-\baselineskip+#1}
  \noindent\rule[0.2ex]{\linewidth}{0.2mm}\par
  \vspace*{\dimexpr-\parskip-.5\baselineskip+#1}}
% Gestrichelte horizontale Linie:
\RequirePackage{dashrule}
\newcommand{\HDRule}[1][\medskipamount]{\par
  \vspace*{\dimexpr-\parskip-\baselineskip+#1}
  \noindent\hdashrule[0.2ex]{\linewidth}{0.2mm}{1mm} \par
  \vspace*{\dimexpr-\parskip-.5\baselineskip+#1}}
% Mathe in Anführungszeichen:
\newsavebox{\mathbox}\newsavebox{\mathquote}
\makeatletter
\newcommand{\mq}[1]{% \mathquotes{<stuff>}
  \savebox{\mathquote}{\text{"}}% Save quotes
  \savebox{\mathbox}{$\displaystyle #1$}% Save <stuff>
  \raisebox{\dimexpr\ht\mathbox-\ht\mathquote\relax}{"}#1\raisebox{\dimexpr\ht\mathbox-\ht\mathquote\relax}{''}
}
\makeatother

% Paragraph mit Zähler (Section-Weise)
\newcounter{cparagraphC}
\newcommand{\cparagraph}[1]{
\stepcounter{cparagraphC}
\paragraph{\thesection{}-\thecparagraphC{} #1}
%\addcontentsline{toc}{subsubsection}{\thesection{}-\thecparagraphC{} #1}
\label{\thesection-\thecparagraphC}
}
\makeatletter
\@addtoreset{cparagraphC}{section}
\makeatother


% (Vorlesungs-)Folien einbinden:
\makeatletter
% Folien-Pfad (inkl. Dateiname ohne Endung und ggf. ohne Nummerierung)
\newcommand{\@slidePath}{}
\newcommand{\slidePath}[1]{\renewcommand{\@slidePath}{#1}}
% Folien Eigenschaften
\newcommand{\@slideScale}{0.5}
\newcommand{\slideScale}[1]{\renewcommand{\@slideScale}{#1}}
\AtBeginDocument{
% Folien von einer Datei skaliert
\newcommand{\slide}[2][\@slideScale]{\slides[#1]{}{#2}}
\newcommand{\slideTrim}[6][\@slideScale]{\slides[#1 , clip,  trim = #5cm #4cm #6cm #3cm]{}{#2}}
% Folien von mehreren nummerierten Dateien skaliert
\newcommand{\slides}[3][\@slideScale]{\begin{center}
\includegraphics[page=#3, scale=#1]{\@slidePath #2.pdf}
\end{center}}
}
\makeatother

% \emph{} anders definieren
\makeatletter
\DeclareRobustCommand{\em}{%
  \@nomath\em \if b\expandafter\@car\f@series\@nil
  \normalfont \else \scshape \fi}
\makeatother

% unwichtiges
\newcommand{\unimptnt}[1]{{\transparent{0.5}#1}}

% alph. enumerate
\newenvironment{anumerate}{\begin{enumerate}[label=(\alph*)]}{\end{enumerate}} % Alphabetische Aufzählung

% Hanging parameters
\newcommand{\hangpara}[1]{\par\noindent\hangindent+2em\hangafter=1 #1\par\noindent}

%% EINFACHE BEFEHLE

% Abkürzungen Mathe
\newcommand{\EE}{\mathbb{E}}
\newcommand{\QQ}{\mathbb{Q}}
\newcommand{\RR}{\mathbb{R}}
\newcommand{\CC}{\mathbb{C}}
\newcommand{\NN}{\mathbb{N}}
\newcommand{\ZZ}{\mathbb{Z}}
\newcommand{\PP}{\mathbb{P}}
\renewcommand{\SS}{\mathbb{S}}
\newcommand{\cA}{\mathcal{A}}
\newcommand{\cB}{\mathcal{B}}
\newcommand{\cC}{\mathcal{C}}
\newcommand{\cD}{\mathcal{D}}
\newcommand{\cE}{\mathcal{E}}
\newcommand{\cF}{\mathcal{F}}
\newcommand{\cG}{\mathcal{G}}
\newcommand{\cH}{\mathcal{H}}
\newcommand{\cI}{\mathcal{I}}
\newcommand{\cJ}{\mathcal{J}}
\newcommand{\cM}{\mathcal{M}}
\newcommand{\cN}{\mathcal{N}}
\newcommand{\cP}{\mathcal{P}}
\newcommand{\cR}{\mathcal{R}}
\newcommand{\cS}{\mathcal{S}}
\newcommand{\cZ}{\mathcal{Z}}
\newcommand{\cL}{\mathcal{L}}
\newcommand{\cT}{\mathcal{T}}
\newcommand{\cU}{\mathcal{U}}
\newcommand{\cX}{\mathcal{X}}
\newcommand{\cV}{\mathcal{V}}
\renewcommand{\phi}{\varphi}
\renewcommand{\epsilon}{\varepsilon}
\renewcommand{\theta}{\vartheta}

% Verschiedene als Mathe-Operatoren
\DeclareMathOperator{\arccot}{arccot}
\DeclareMathOperator{\arccosh}{arccosh}
\DeclareMathOperator{\arcsinh}{arcsinh}
\DeclareMathOperator{\arctanh}{arctanh}
\DeclareMathOperator{\arccoth}{arccoth} 
\DeclareMathOperator{\var}{Var} % Varianz 
\DeclareMathOperator{\cov}{Cov} % Co-Varianz 

% Farbdefinitionen
\definecolor{red}{RGB}{180,0,0}
\definecolor{green}{RGB}{75,160,0}
\definecolor{blue}{RGB}{0,75,200}
\definecolor{orange}{RGB}{255,128,0}
\definecolor{yellow}{RGB}{255,245,0}
\definecolor{purple}{RGB}{75,0,160}
\definecolor{cyan}{RGB}{0,160,160}
\definecolor{brown}{RGB}{120,60,10}

\definecolor{itteny}{RGB}{244,229,0}
\definecolor{ittenyo}{RGB}{253,198,11}
\definecolor{itteno}{RGB}{241,142,28}
\definecolor{ittenor}{RGB}{234,98,31}
\definecolor{ittenr}{RGB}{227,35,34}
\definecolor{ittenrp}{RGB}{196,3,125}
\definecolor{ittenp}{RGB}{109,57,139}
\definecolor{ittenpb}{RGB}{68,78,153}
\definecolor{ittenb}{RGB}{42,113,176}
\definecolor{ittenbg}{RGB}{6,150,187}
\definecolor{itteng}{RGB}{0,142,91}
\definecolor{ittengy}{RGB}{140,187,38}

% Textfarbe ändern
\newcommand{\tred}[1]{\textcolor{red}{#1}}
\newcommand{\tgreen}[1]{\textcolor{green}{#1}}
\newcommand{\tblue}[1]{\textcolor{blue}{#1}}
\newcommand{\torange}[1]{\textcolor{orange}{#1}}
\newcommand{\tyellow}[1]{\textcolor{yellow}{#1}}
\newcommand{\tpurple}[1]{\textcolor{purple}{#1}}
\newcommand{\tcyan}[1]{\textcolor{cyan}{#1}}
\newcommand{\tbrown}[1]{\textcolor{brown}{#1}}

% Umstellen der Tabellen Definition
\newcommand{\mpb}[1][.3]{\begin{minipage}{#1\textwidth}\vspace*{3pt}}
\newcommand{\mpe}{\vspace*{3pt}\end{minipage}}

\newcommand{\resultul}[1]{\underline{\underline{#1}}}
\newcommand{\parskp}{$ $\\}	% new line after paragraph
\newcommand{\parnl}{\unskip\parfillskip 0pt \par} % justified new line (brace line!)

% mathematische Abkürzungen
\newcommand{\corr}{\;\widehat{=}\;}
\newcommand{\mdeg}{^{\circ}}
\newcommand{\nok}[2]{\binom{#1}{#2}}	% n über k BESSER: \binom{n}{k}
\newcommand{\mtr}[1]{\begin{pmatrix}#1\end{pmatrix}}	% Matrix
\newcommand{\dtr}[1]{\begin{vmatrix}#1\end{vmatrix}}	% Determinante (Betragsmatrix)
\LetLtxMacro{\originalVec}{\vec}
\renewcommand{\vec}[1]{\underline{#1}}	% Vektorschreibweise
\newcommand{\imptnt}[1]{\colorbox{red!30}{#1}}	% Wichtiges
\newcommand{\intd}[1]{\,\mathrm{d}#1}
\newcommand{\diffd}[1]{\mathrm{d}#1}
% für Module-Rechnung: \pmod{}
\newcommand{\unit}[1]{\,\mathrm{#1}}
%\LetLtxMacro{\ntilde}{\tilde}
%\renewcommand{\tilde}{\widetilde}

% weiter Abkürzungen
\newcommand{\gdw}{genau dann wenn}
\newcommand{\lecdate}[1]{\marginpar{\small\textcolor{gray}{Vorlesung\newline #1}}}
\definecolor{htworange}{RGB}{249,155,28}	% Corporate Design orange CMYK(0/45/100/0)
\definecolor{htwblue}{RGB}{0,116,188}			% Corporate Design blau CMYK(90/50/0/0)
\definecolor{htwgrey}{RGB}{128,128,128}		% Corporate Design gray 60% Schwarz [Annahme: CMYK(0,0,0,60)]
\RequirePackage{helvet}	% Corporate Design schreibt Arial/Helvetica vor
\renewcommand{\familydefault}{\sfdefault}	% s.o.

\usenavigationsymbolstemplate{}
\setbeamercolor{author in head/foot}{fg=black}
\setbeamercolor{title}{fg=black}
\setbeamercolor{bibliography entry author}{fg=htworange!70}
%\setbeamercolor{bibliography entry title}{fg=blue} 
\setbeamercolor{bibliography entry location}{fg=htworange!60} 
\setbeamercolor{bibliography entry note}{fg=htworange!60}  
\setbeamercolor{caption name}{fg=htworange}
\setbeamercolor{frametitle}{fg=black}
\setbeamerfont{frametitle}{size=\large}
\setbeamertemplate{frametitle}{\hspace*{-.75em}\insertframetitle}

\setbeamercolor{itemize item}{fg=black,bg=white}
\setbeamercolor{enumerate item}{fg=black,bg=white}
\setbeamercolor{itemize subitem}{fg=black,bg=white}
\setbeamercolor{enumerate subitem}{fg=black,bg=white}
\setbeamercolor{itemize subsubitem}{fg=black,bg=white}
\setbeamercolor{enumerate subsubitem}{fg=black,bg=white}
\setbeamertemplate{itemize item}{$\bullet$}
\setbeamertemplate{itemize subitem}{--}
\setbeamertemplate{itemize subsubitem}{$\bullet$}
\makeatletter
\setbeamertemplate{footline}
{
\leavevmode
\def\arraystretch{1.2}
\arrayrulecolor{gray}
\begin{tabular}{ p{0.167\textwidth} | p{0.491\textwidth} | p{0.089\textwidth} | p{0.103\textwidth}}
\hline
\strut\insertshortauthor & \insertshorttitle & Folie \insertframenumber{}% / \inserttotalframenumber{}
 & \@date\\
\end{tabular}
}
\setbeamertemplate{headline}
{%
\leavevmode%
\setlength{\arrayrulewidth}{1pt}%
\hspace*{2em}	%
\begin{tabular}{p{0.63\textwidth}}%
\rule{0pt}{3em}{\normalsize%
\textbf{%
\insertsectionhead%
%\setbox0=\hbox{\subsecname\unskip}\ifdim\wd0=0pt\else%
%~--~\insertsubsectionhead%
%\fi%
\ifx\insertsubsection\empty\else%
 \ifdefempty{\subsecname}{\relax}{%
  ~--~\insertsubsectionhead%
 }%
\fi
\strut}%
}\\%
\arrayrulecolor{htworange}%
\hline%
\end{tabular}%
\begin{tabular}{l}%
\rule{0pt}{4em}%
\includegraphics[width=3.2cm]{\pathtomaster/HTW.eps}\\%
\end{tabular}%
}%
\makeatother

% Recalculate the head of head- and footlines for each frame
% (http://tex.stackexchange.com/a/61788/3323)
\usepackage{etoolbox}
\makeatletter
\pretocmd{\beamer@@@@frame}{%
  \addtocounter{framenumber}{1}%
  \beamer@calculateheadfoot%
  \addtocounter{framenumber}{-1}}{}{}
\makeatother

\author{Falk-Jonatan Strube}
\title{\texttt{ledleuchtturm} -- Dokumentation}
\date{18.07.2017}

\usepackage{pgfpages}
\setbeameroption{show notes on second screen=right}
\begin{document}
\author{Falk-Jonatan Strube}
\title{\texttt{ledleuchtturm} -- Dokumentation}

\begin{frame}[plain]
\begin{center}
\includegraphics[width=.7\textwidth]{\pathtomaster/HTW.eps}
\end{center}
\maketitle
\end{frame}

\section{Gliederung}
\begin{frame}
\setstretch{1.5}
\note{\setstretch{1.5}}
\begin{itemize}
\item<2-> Aufgaben
\begin{itemize}
\item<3-> Umfang
\item<4-> Durchführung
\end{itemize}
\item<5-> Bewertung
\begin{itemize}
\item<6-> Aufgaben
\note<6>[item]{Wie ist die Aufgabe zu bewerten?}
\item<7> Projekt
\note<7>[item]{NEXT}
\end{itemize}
\end{itemize}
\end{frame}

\section{Aufgaben}
\subsection{Umfang}
\begin{frame}
\setstretch{1.5}
\note{\setstretch{1.5}}
\note{Aufgaben lassen sich in unterschiedliche Kategorien packen:}
\begin{itemize}
\item<2-> Als \emph{Dokument}\,ierender
\item<3-> Delegierte Aufgaben
\item<4-> Weitere
\note<4>[item]{NEXT}
\end{itemize}
\end{frame}

\begin{frame}
\setstretch{1.5}
\note{\setstretch{1.5}}
Als \emph{Dokument}\,ierender:
\begin{itemize}
\item<2-> \LaTeX{}-Vorlagen
\note<2>[item]{Vorlagen im Corporate Design der HTW (Stil)}
\note<2>[item]{Hilfestellung}
\item<3-> Dokumentenvorlagen zum „Ausfüllen“
\note<3>[item]{Struktur}
\note<3>[item]{inhaltliche Vorüberlegungen.}
\item<4-> Code-Dokumentation (Epydoc)
\note<4>[item]{Beispiel erstellt}
\note<4>[item]{Hilfestellung und Rückfragen}
\note<4>[item]{Anpassung an \LaTeX{} (Stil, Übersetzung).}
\item<5-> Koordination Dokumentationsarbeit
\note<5>[item]{Wer muss was schreiben?}
\item<6-> Überarbeitung der Dokumente
\note<6>[item]{Korrekturen}
\note<6>[item]{teils große Änderungen}
\note<6>[item]{Übersetzung ins Deutsche}
\note<6>[item]{NEXT}
\end{itemize}
\end{frame}

\begin{frame}
\setstretch{1.5}
\note{\setstretch{1.5}}
Delegierte Aufgaben:
\begin{itemize}
\item<2-> Pflichtenheft (Überarbeitung)
\note<2>[item]{Finale Überarbeitung}
\note<2>[item]{Ergänzte Inhalte (Abrundung)}
\item<3-> Vorlage Unit-Tests
\note<3>[item]{Erarbeiten, wie Unit Tests in Python funktionieren.}
\item<4-> Vorlage Python Paketstruktur
\note<4>[item]{Von Python Skriptsammlung $\Rightarrow$ gebündelte Paketstruktur}
\item<5-> \lstinline`changeState()`-Funktionen
\note<5>[item]{Module LED-Funktionen}
\end{itemize}
\uncover<6->{Weitere:}
\begin{itemize}
\item<7-> \lstinline`git`
\note<7>[item]{Schon in SE1 angedeutet}
\note<7>[item]{Rückfragen}
\note<7>[item]{Grundprinzip}
\note<7>[item]{Einrichtung}
\item<8-> Prototyp MQTT (verworfen)
\note<8>[item]{Prototyp entwickelt}
\note<8>[item]{Vor- und Nachteile erarbeitet.}
\note<8>[item]{NEXT: Nun an einigen Beispielen Durchführung erläutern.}
\end{itemize}
\end{frame}

\subsection{Durchführung}
\begin{frame}{\LaTeX{}-Vorlage (Layout)}
\setstretch{1.5}
\note{\setstretch{1.5}}
Styleguide der HTW Dresden:
\note<1>[item]{Danach gerichtet, um Vorlage zu erstellen}
\begin{itemize}
\item<2-> Schrift
\note<2>[item]{beinhaltet:}
\begin{itemize}
\item<3-> Schriftart
\note<3>[item]{Arial/Helvetica}
\item<4-> Schriftgröße
\end{itemize}
\item<5-> Logo
\note<5>[item]{Dazu gab es:}
\begin{itemize}
\item<6-> Positionierung
\item<7-> Größe
\end{itemize}
\item<8-> Vorgaben für wissenschaftliche Dokumente?!
\note<8>[item]{Keine weiteren Vorgaben vorhanden. vgl. TU Dresden}
\note<8>[item]{NEXT: Beispiel}
\end{itemize}
\end{frame}

\begin{frame}{\LaTeX{}-Vorlage (Layout)}
\setstretch{1.5}
\note{\setstretch{1.5}}
%trim: {left lower right upper}
\begin{center}
\includegraphics[scale=.27, page=1, trim={2cm 2cm 2cm 2cm},clip]{Beispiel_Layout.pdf}
\includegraphics[scale=.27, page=2, trim={2cm 2cm 2cm 2cm},clip]{Beispiel_Layout.pdf}
\end{center}
\note{
\begin{itemize}
\item Orientiert an CD der TUD:
\begin{itemize}
\item Aufmachung der Titelseite(n)
\item Layout der Überschriften
\end{itemize}
\item Einhaltung des CD der HTW
\begin{itemize}
\item Logo
\item Schriftart
\end{itemize}
\end{itemize}
}
\end{frame}

\begin{frame}{\LaTeX{}-Vorlage (Struktur)}
\setstretch{1.5}
\note{\setstretch{1.5}}
\note<1>[item]{Struktur!!!}
\begin{itemize}
\item<2-> Benötigte Dokumenteninhalte
\note<2>[item]{Was gehört in \\
Protokolle\\
Entwicklerdoku\\
Projektdoku\\
Gesamtdoku\\
Präsentation?}
\item<3-> Optimale Struktur
\note<3>[item]{Sinnvolle Reihenfolge?}
\note<3>[item]{NEXT: Beispiel}
\end{itemize}
\end{frame}

\begin{frame}[fragile]{\LaTeX{}-Vorlage (Struktur)}
%\setstretch{1.5}
\note{\setstretch{1.5}}
\begin{lstlisting}[language={[LaTeX]TeX}, basicstyle=\ttfamily\footnotesize, frame=none, numbers=none]
\maketitle
\section{Gruppentreffen ???}
% Tabellarische Auflistung von Beginn, Ende, Ort, Anwenden und Protokollant

\paragraph{Tagesordnungspunkte}
% Auflistung der TOPs

\subsection{TOP ???}
% TOPs als Unterabschnitte

\subsection{Nächste Schritte}
% Tabellarische Auflistung der nächsten Schritte
\end{lstlisting}
\note{
Protokoll:
\begin{itemize}
\item Benötigte Dokumenteninhalte, nicht nur für Protokoll!\\
Was gehört in Protokolle, Entwickler-, Projekt-, Gesamtdoku und Präsentation?
\item Optimale Struktur\\
Sinnvolle Reihenfolge?
\end{itemize}
}
\end{frame}

%\begin{frame}[fragile]{\LaTeX{}-Vorlage (Struktur)}
%%\setstretch{1.5}
%\note{\setstretch{1.5}}
%\begin{lstlisting}[language={[LaTeX]TeX}, basicstyle=\ttfamily\footnotesize, frame=none, numbers=none]
%% Entwicklerdokumentation = Programmiererdokumentation (Pflichtenheft /DA2/)
%\section{Grobentwurf}
%% Grobentwurf aus Pflichtenheft entnehmen und noch mal als Dokumentation aufarbeiten.
%\section{Feinentwurf}
%% Feinentwurf aus Entwurf/Entwurf.tex noch mal aufarbeiten und Dokumentieren.
%\section{Struktur des Codes}
%% Struktur des Codes (wie er abgegeben wurde) festhalten. 
%% Wie sind die Dateien/Packet(e) organisiert?
%% Alle Dateien beschreiben: Überblick verschaffen.
%% Am besten (statt Klassendiagramm) ein "Datei-Diagramm" erstellen.
%% [...]
%\end{lstlisting}
%\note{Entwicklerdokumentation:
%\begin{itemize}
%\item Benötigte Dokumenteninhalte\\
%Was gehört in Protokolle, Entwickler-, Projekt-, Gesamtdoku und Präsentation?
%\item Optimale Struktur\\
%Sinnvolle Reihenfolge?
%\end{itemize}
%}
%\end{frame}

\begin{frame}{Pflichtenheft}
\setstretch{1.5}
\note{\setstretch{1.5}}
\begin{itemize}
\item<2-> Inhaltliche Verbesserungen
\item<3-> Ergänzungen
\item<4-> Verschönerung
\item<5-> Strukturierung
\note<5>[item]{NEXT: Beispiel}
\end{itemize}
\end{frame}

\begin{frame}{Pflichtenheft (Entwurf $\Rightarrow$ Abgabedokument)}
\setstretch{1.5}
\note{\setstretch{1.5}}
%trim: {left lower right upper}
\begin{center}
\includegraphics[scale=.27, page=1, trim={2cm 2cm 2cm 2cm},clip]{Beispiel_Pflichtenheft_first.pdf}
\includegraphics[scale=.27, page=1, trim={2cm 2cm 2cm 2cm},clip]{Beispiel_Pflichtenheft_final.pdf}
\end{center}
\note{
\begin{itemize}
\item Inhaltliche Verbesserungen\\
Mehrdeutigkeiten, Redundanz (Qualitätsanf.)
\item Ergänzungen\\
Eingangstext, Beschreibung
\item Verschönerung\\
Verdeutlichung der Zustände, Unterschriftzeile
\item Strukturierung\\
Nummerierung der Anforderungen zum referenzieren
\end{itemize}
}
\end{frame}

%\begin{frame}{EpyDoc}
%\setstretch{1.5}
%\note{\setstretch{1.5}}
%\note[item]{Ähnlich wie JavaDoc}
%\begin{itemize}
%\item Kommentierung des Codes
%\item Generierung einer Dokumentation
%\begin{itemize}
%\item \LaTeX{}-Ausgabe
%\item Anpassungen
%\item Übersetzung
%\end{itemize}
%\end{itemize}
%\end{frame}
%
%\begin{frame}[fragile]{EpyDoc}
%%\setstretch{1.5}
%\note{\setstretch{1.5}}
%\begin{lstlisting}[language={Python}, basicstyle=\ttfamily\footnotesize, frame=none, numbers=none]
%def testFunc(test):
%    '''
%    Beschreibung
%    
%    @type test: Typ Parameter
%    @param test: Beschreibung Parameter
%    @rtype: Typ Rückgabewert
%    @return: Beschreibung Rückgabewert
%    @author: Autor
%    '''
%\end{lstlisting}
%\end{frame}
%
%\begin{frame}[fragile]{EpyDoc}
%%\setstretch{1.5}
%\note{\setstretch{1.5}}
%%trim: {left lower right upper}
%\begin{center}
%\includegraphics[scale=.5, page=1, trim={3cm 15.5cm 3cm 7.5cm},clip]{Beispiel_EpyDoc.pdf}
%\end{center}
%\end{frame}

\begin{frame}{LED-Funktionen}
\setstretch{1.5}
\note{\setstretch{1.5}}
Funktion \lstinline`changeState()`
\note<1>[item]{Je nach Status: LED an-/ausschalten oder blinken. Dafür benötigte Funktionen:}
\begin{itemize}
\item<2-> \lstinline`toggleLED()`, \lstinline`enableLED()`
\note<2>[item]{Einfache Zustandsänderung}
\item<3-> \lstinline`switchAllOff()`
\note<3>[item]{Zurücksetzen}
\item<4-> Blinken:
\begin{itemize}
\item<5-> \lstinline`blink()`
\item<6-> \lstinline`blinkThreadFunction()`
\note<6>[item]{Thread als Funktion}
\item<7-> \lstinline`killBlinkThread()`
\note<7>[item]{Thread wird nur gestartet, wenn nötig und sofort wieder beendet, wenn nicht mehr benötigt)}
\note<7>[item]{NEXT}
\end{itemize}
\end{itemize}
\end{frame}

\section{Bewertung}
\subsection{Aufgaben}
\begin{frame}
\setstretch{1.5}
\note{\setstretch{1.5}}
\begin{itemize}
\item<2-> Python
\note<2>[item]{Wissen erweitert}
\note<2>[item]{Anwendung von Threads: BS2, PKA}
\item<3-> \LaTeX{}
\note<3>[item]{Wissen erweitert}
\item<4-> \texttt{git}
\note<4>[item]{Wissen erweitert}
\note<4>[item]{Jetzt auch in Gruppe und nicht nur alleine genutzt}
\item<5-> Dokumentieren
\note<5>[item]{Erfahrung gesammelt}
\note<5>[item]{Strukturierung von Doku}
\note<5>[item]{NEXT}
\end{itemize}
\end{frame}

\subsection{Projekt}
\begin{frame}
\setstretch{1.5}
\note{\setstretch{1.5}}
\begin{itemize}
\item<2-> Gruppengröße 
\note<2-3>[item]{Gruppe zu groß}
\item<3-> Umfang der Aufgabe 
\note<3>[item]{Aufgabe zu klein: geht nicht größer, weil Belege}
\item<4-> Unterschiedliche Motivation / Anspruch\\
$\Rightarrow$ Entmutigt Delegation
\note<4>[item]{Unterschiedliche Motivation/Anspruch \\
$\to$ mehr Arbeit für Gruppenmitglieder mit höherem Anspruch}
\item<5-> Doch Objektorientierung?
\note<5>[item]{aktuelle Lösung gut}
\note<5>[item]{Objektorientierung:\\
würde einiges einfacher machen \\
haben schon Objekte (Thread, …)}
\item<6-> Lieber MQTT?
\note<6>[item]{MQTT: Sockets zwar nicht ohne Grund gewählt (Broker)}
\note<6>[item]{MQTT allerdings dynamischer}
\note<6>[item]{DONE}
\end{itemize}
\end{frame}

\end{document}